\def\topic{}
\ifdefined\MyHeaderIncluded
\else
\def\qwertyflag{true}

\ifdefined\MyHeaderIncluded

\else
  \newcommand{\foobar}{FooBar}
  \def\MyHeaderIncluded{true}
% if is till the end of the file

\documentclass[10pt,a4paper]{article}
\usepackage[utf8]{inputenc}
%\usepackage[german]{babel}
%\usepackage[cyrillic]{babel}
\usepackage[T1]{fontenc}
%\usepackage[T2A]{fontenc}
%\usepackage[english, russian]{babel}
\usepackage{amsmath}
\usepackage{amsfonts}
\usepackage{amssymb}
\usepackage{amsthm}
\usepackage{mathtools} %coloneqq

\usepackage{bussproofs} %

%\renewcommand{\thesubsection}{\arabic{subsection}}
%\newtheorem{thm}{Theorem}[section]
%\renewcommand{\thethm}{\arabic{subsection}.\arabic{thm}}

\usepackage{thmtools}
\usepackage{hyperref}
\hypersetup{
    colorlinks=true,
    linkcolor=blue,
    filecolor=magenta,      
    urlcolor=cyan,
    pdftitle={Overleaf Example},
    %pdfpagemode=FullScreen,
    }
%\usepackage{stix-type1}%\turnediota
\theoremstyle{definition}

\newtheorem{constructor}{Constructor}[section]
\newtheorem{iaxiom}{Axiom}
\newtheorem{axiom}{Axiom}
\newtheorem*{axiom*}{Axiom}
\newtheorem{definition}{Definition}[section]
\newtheorem*{definition*}{Definition}
\newtheorem{remark}{Remark}[section]
\newtheorem*{remark*}{Remark}
\newtheorem{recall}{Recall}[section]
\newtheorem*{recall*}{Recall}

\newtheorem{lemma}{Lemma}[section]
\newtheorem{theorem}{Theorem}[section]
\newtheorem*{theorem*}{Theorem}
\newtheorem{note}{Note}[section]
\usepackage{etoolbox}

%\newcommand*{\@Dashv}[2]{%
%  \reflectbox{$\m@th#1#2$}%
%}
\newcommand*{\separation}[3]{% 1)binder, 2)set, 3)condition
  \{#1\in #2: #3\}
}

% https://www.freecodecamp.org/news/simple-regex-tricks-for-beginners-3acb3fa257cb/
%  \\mathcal\{T\} |---> \topology
%  \\mathcal\{([^\}]*)\}   |--->  $1
\newcommand{\AND}{\ \ \&\ \ }
\newcommand{\any}{\mathrm{\ any}}
\newcommand{\borel}{\mathcal{B}}
\newcommand{\can}{\mathrm{can}}
\newcommand{\canonize}{\mathrm{canonize}}
\newcommand{\closure}{\mbox{cl}}
\newcommand{\complex}{\mathbb{C}}
\newcommand{\converges}{\longrightarrow}
\newcommand{\countable}{\mbox{\ countable}}
\newcommand{\ldefined}{\mbox{\ defined}}
\newcommand{\dom}{\mbox{dom}}
\newcommand{\Dom}{\mbox{Dom}}
\newcommand{\Done}{\noindent \textbf{Done.}\\}
\newcommand{\lexists}{\mbox{\ exists}}
\newcommand{\lforall}{\mbox{\ forall}}
\newcommand{\lunique}{\mbox{\ unique}}
\newcommand{\family}{\mathrm{\ family}}
\newcommand{\Fm}{\mathrm{Fm}}
\newcommand{\FuncSymb}{\mathrm{FuncSymb}}
\newcommand{\function}{\mathrm{\ function}}
\newcommand{\functional}{\mbox{\ functional}}
\newcommand{\IFF}{\ \Longleftrightarrow\ }
\renewcommand{\iff}{\ \longleftrightarrow\ }
\newcommand{\IMPLIES}{\ \Longrightarrow\ }
\renewcommand{\implies}{\ \longrightarrow\ }
\newcommand{\Inf}{\mathrm{Inf}}
\newcommand{\injective}{\mbox{\ injective}}
\newcommand{\injection}{\mbox{\ injection}}
\newcommand{\interior}{\mbox{int}}
\newcommand{\intersection}{\bigcap}
\newcommand{\isempty}{\mathrm{\ empty}}
\newcommand{\LB}{\mathrm{LB}}
\newcommand{\nghbset}{\mathcal{N}}
\newcommand{\nonempty}{\mathrm{\ nonempty}}
\newcommand{\Ord}{\mathrm{Ord}}
\newcommand{\power}{\mathcal{P}}
\newcommand{\pr}{\mathrm{pr}}
\newcommand{\fst}{\mathrm{fst}}
\newcommand{\snd}{\mathrm{snd}}
\newcommand{\PredSymb}{\mathrm{PredSymb}}
\newcommand{\Proof}{\noindent \textbf{Proof.}\\}
\newcommand{\ran}{\mathrm{ran}}
\newcommand{\Ran}{\mathrm{Ran}}
\newcommand{\real}{\mathbb{R}}
\newcommand{\relation}{\mathrm{\ relation}}
\newcommand{\set}{\mathrm{\ set}}
\newcommand{\singleton}{\mathrm{\ singleton}}
\newcommand{\semantics}{\mathrm{semantics}}
\newcommand{\semialgebra}{\mathcal{S}}
\newcommand{\Sup}{\mathrm{Sup}}
\newcommand{\surjective}{\mbox{\ surjective\ on\ }}
\newcommand{\Term}{\mathrm{Term}}
\newcommand{\Type}{\mathrm{Type}}
\newcommand{\topology}{\mathcal{T}}
\newcommand{\total}{\mbox{\ total\ on\ }}
\newcommand{\transitive}{\mathrm{\ transitive}}
\newcommand{\union}{\bigcup}
\newcommand{\lTrue}{\mathrm{True}}
\newcommand{\lFalse}{\mathrm{False}}
\newcommand{\lUncertain}{\mathrm{Uncertain}}
\newcommand{\lif}{\mathrm{if}}
\newcommand{\lthen}{\mathrm{then}}
\newcommand{\lelse}{\mathrm{else}}
\newcommand{\lexception}{\mathrm{exception}}
\newcommand{\lundef}{\mathrm{undef}}
\newcommand{\univocal}{\mathrm{\ univocal}}
\newcommand{\univocalat}{\mathrm{\ univocal\ at\ }}
\newcommand{\UB}{\mathrm{UB}}
\newcommand{\Un}{\mathcal{U}}
\newcommand{\Var}{\mathrm{Var}}

\newcommand{\notvarphi}{\phi}
\renewcommand{\phi}{\varphi}

\newcommand{\comment}[1]{\quad\mbox{(#1)}}
% local commands
%\preto\section{\undef\blipo} %etoolbox
\preto\subsection{\undef\blipo} %etoolbox
% https://tex.stackexchange.com/questions/105684/limit-macro-newcommand-scope-to-section
\usepackage{graphicx}
\usepackage{authblk}

\setlength\parindent{0pt} % whole doc of \noindent
\usepackage{soul}
%\usepackage{mathgl}

\fi % guard
\ifdefined\topic
\else
\def\topic{No Titel}
\fi
\title{\topic}
\author[1]{Georgy Dunaev}
\begin{document}
\maketitle
\fi
\section{Proving limitation of size}
\subsection{Foundations and Large Cardinals}%
This part is self-contained.\\
The goal is to prove axiom of limitation of size using small amount of simple axioms and definitions.\\
\begin{iaxiom}
Axiom of limitation of size:
\[
(K\in V)
\iff
(|K| < |V|)
\]
\end{iaxiom}


Here $|\cdot|$ is a functional symbol defined as following.
\begin{definition}(cardinality of a class)
\[
|K|:= \Ord \cap \bigcap \{a\in\Ord|K\preccurlyeq a\}
\]
\end{definition}
\begin{definition}
\[
(X\preccurlyeq Y)
\iff
(\exists f: (f:X\to Y)\land (f \injective))
\]
\end{definition}
The intended behaviour of this definition is to have cardinality of argument if the argument of a functional symbol is a set, and $\Ord$ if the argument is a proper class.\\
We will introduce some axioms to make such behaviour provable.\\
\begin{iaxiom}
\[
|\cdot|:V \to \Ord
\]
Actually, since $|\cdot|$ is a functional symbol, we mean the following:
\[
(x\mapsto |x|):V \to \Ord
\]
Functional symbol is defined on any class as argument, while $(x\mapsto |x|)$ is defined on sets only. By definition, $(x\mapsto |x|)$ is a class function defined as a class of pairs $\{\langle x, |x|\rangle | x \any\set\}$.
\end{iaxiom}
Actually, we can use the following axiom:
\begin{axiom*}
\[
(x\set)\implies(|x|\in\Ord)
\]
\end{axiom*}
or even
\begin{axiom*}
\[
(x\set)\implies(|x|\set)
\]
\end{axiom*}
since later we will prove that
\begin{theorem*}
\[
(x\set)\implies(|x|\in\Ord)
\]
\end{theorem*}

\begin{iaxiom} Sets are hereditary:
\[
(A\subseteq B)\land (B\set) \implies (A\set)
\]
\end{iaxiom}

Some equivalent(in normal context) statements:
\[
(K\in V)
\iff
(|K| \subsetneq |V|)
\]
\[
(K\in V)
\iff
(|K| \neq |V|)
\]
because $|V|=\Ord$ and
\begin{theorem}
\[
|V|=\Ord
\]
\end{theorem}
\Proof
\Done

In this text we will consider the simplest definition of ordinals which assumes that axiom of regularity holds. It is done for shortness, and can be done with a definition that doesn't require the regularity axiom.
\begin{definition} predicates on classes "is transitive" and "is ordinal"
\[
tr(X) \iff (\forall Z\in X: Z\subseteq X)
\]
\[
ord(X) \iff tr(X)\land (\forall Y\in X: tr(Y))
\]
the other possibilities are
\[
ord(X) \iff tr(X)\land (\mbox{E Wo }X)
\]
\[
ord(X) \iff tr(X)\land (\mbox{E Wo }X)\land (\mbox{E Fr }X)
\]
etc ...
\end{definition}
\[
(X \set)
\implies
\]

\subsection{Proof of the main statement}
We want to prove the following statement:
\[
(K \in V) \iff (|K| \neq |V|)
\]

Firstly, we will prove some lemmas:
\begin{lemma}
\[
\Ord\notin\Ord
\]
\end{lemma}
\Proof
Assume $\Ord\in\Ord$.\\
Therefore $\in$ well-orders $\Ord$ as a well-founded strict linear order.\\
In particular strict order is irreflexive.\\
Since using assumpton, $\Ord$ is an element of $\Ord$, by irreflexetivity, we have $\Ord\notin\Ord$.
Contradiction with assumption.
Therefore $\Ord\notin\Ord$ is proven.
\Done

\begin{lemma}
\[
|V|=\Ord
\]
\end{lemma}
\Proof
...\\
\Done

\begin{lemma}(1st part of the statement, from left to right)
\[
(K \in V) \implies (|K| \neq |V|)
\]
\end{lemma}
\Proof
Assume $(K \in V)$.\\
By our axiom, since $(K \in V)$ we have $(|K| \in \Ord)$.\\
By previous lemmas we have $|V|=Ord$ and $\Ord\notin\Ord$.\\
From these three statements we deduce that $|K| \neq |V|$.\\
\Done

\begin{lemma}(2nd part of the statement, from right to left)
\[
(|K| \neq |V|) \implies (K \in V)
\]
\end{lemma}
\Proof
Assume $(|K| \neq |V|)$.\\
We know that $(|V|=\Ord)$, so $(|K| \neq Ord)$.\\
Recalling the definition:
$|K|:= \Ord \cap \bigcap \{a\in\Ord|K\preccurlyeq a\}$.\\
Now we have that 
there exists some $m \in \Ord$, which is not in $\bigcap \{a\in\Ord|K\preccurlyeq a\}$. (vice versa situation is not possible since any non-emptyintersection of subclasses of $\Ord$ is subclass of $\Ord$)\\
...

\Done

\ifdefined\qwertyflag
\undef\qwertyflag
\end{document}
\fi